% from ega0-1-fr.tex
\oldpage[0\textsubscript{I}]{22}

\begin{env}[2.1.5]
\label{0.2.1.5-fr}
Soient $X$ un espace irréductible, $f$ une application continue de $X$ dans un
espace topologique $Y$. Alors $f(X)$ est irréductible, et si $x$ est point
générique de $X$, $f(x)$ est point générique de $f(X)$ et par suite aussi de
$\overline{f(X)}$. En particulier, si de plus $Y$ est irréductible et a un seul
point générique $y$, pour que $f(X)$ soit partout dense, il faut et il suffit
que $f(x)=y$.
\end{env}

\begin{env}[2.1.6]
\label{0.2.1.6-fr}
Tout sous-espace irréductible d'un espace topologique $X$ est contenu dans un
sous-espace irréductible \emph{maximal}, qui est nécessairement fermé. Les
sous-espaces irréductibles maximaux de $X$ sont appelés les \emph{composantes
irréductibles} de $X$. Si $Z_1$, $Z_2$ sont deux composantes irréductibles
distinctes de l'espace $X$, $Z_1\cap Z_2$ est un ensemble fermé \emph{rare}
dans chacun des sous-espaces $Z_1$, $Z_2$; en particulier, si une composante
irréductible de $X$ admet un point générique (2.1.2) un tel point ne peut
appartenir à aucune autre composante irréductible. Si $X$ n'a qu'un nombre
\emph{fini} de composantes irréductibles $Z_i$ ($1\leq i\leq n$), et si, pour
chaque $i$, on pose
\[
  U_i=Z_i\cap\complement\Bigl(\bigcup_{j\neq i}Z_j\Bigr),
\]
les $U_i$ sont ouverts, irréductibles, deux à deux sans point commun et leur
réunion est dense dans $X$.

Soit $U$ une partie ouverte d'un espace topologique $X$. Si $Z$ est une partie
irréductible de $X$ rencontrant $U$, $Z\cap U$ est ouvert et dense dans $Z$,
donc irréductible; inversement, pour toute partie fermée irréductible $Y$ de
$U$, l'adhérence $\overline Y$ de $Y$ dans $X$ est irréductible et
$\overline Y\cap U=Y$. On en conclut qu'il y a \emph{correspondance
biunivoque} entre les composantes irréductibles de $U$ et les composantes
irréductibles de $X$ qui rencontrent $U$.
\end{env}

\begin{env}[2.1.7]
\label{0.2.1.7-fr}
Si un espace topologique $X$ est réunion d'un nombre \emph{fini} de
sous-espaces irréductibles fermés $Y_i$, les composantes irréductibles de $X$
sont les éléments maximaux de l'ensemble des $Y_i$, car si $Z$ est une partie
fermée irréductible de $X$, $Z$ est réunion des $Z\cap Y_i$, d'où on tire
aussitôt que $Z$ doit être contenu dans un des $Y_i$. Soit $Y$ un sous-espace
d'un espace topologique $X$, et supposons que $Y$ n'ait qu'un nombre fini de
composantes irréductibles $Y_i$ ($1\leq i\leq n$); alors les adhérences
$\overline{Y_i}$ dans $X$ sont les composantes irréductibles de $\overline Y$.
\end{env}

\begin{env}[2.1.8]
\label{0.2.1.8-fr}
Soit $Y$ un espace irréductible admettant un seul point générique $y$. Soient
$X$ un espace topologique, $f$ une application continue de $X$ dans $Y$.
Alors, pour toute composante irréductible $Z$ de $X$ rencontrant $f^{-1}(y)$,
$f(Z)$ est dense dans $Y$. La réciproque n'est pas nécessairement vraie;
toutefois, si $Z$ admet un point générique $z$, et si $f(Z)$ est dense dans
$Y$, on a nécessairement $f(z)=y$ (2.1.5); en outre, $Z\cap f^{-1}(y)$ est
alors l'adhérence de $\{z\}$ dans $f^{-1}(y)$ et est donc irréductible, et
comme toute partie irréductible de $f^{-1}(y)$ contenant $z$ est
nécessairement contenue dans $Z$ (2.1.6), $z$ est point générique de
$Z\cap f^{-1}(y)$. Comme toute composante irréductible de $f^{-1}(y)$ est
contenue dans une composante irréductible de $X$, on voit que si toute
composante irréductible $Z$ de $X$ rencontrant $f^{-1}(y)$ admet un point
générique, alors il y a \emph{correspondance biunivoque} entre l'ensemble de
ces composantes et l'ensemble des composantes irréductibles
$Z\cap f^{-1}(y)$ de $f^{-1}(y)$, les points génériques de $Z$ étant
identiques à ceux de $Z\cap f^{-1}(y)$.
\end{env}

